\scanpage{209}\mspage{- 51 -}
\hfill IX

5. \uline{Le théorème d'orthogonalité dans le cas} $\handwritten{\ell}=p$ (\uline{cas semi-stable}). \uline{Dualité des schémas abéliens} $B_{\unsure{\mathrm o|0}}$ \uline{et} $B'_{\mathrm o}$. \uline{Caractérisation de la partie fixe}. \uline{Critères de bonne réduction}, \uline{de bonne réduction essentielle et semi-stable}.

5.1. Plaçons\unsure{-|\,}nous dans le cas 2.2.3 où $S$ est hensélien et $A^{\mathrm o}$ est $\handwritten{\ell}$-divisible, de sorte que la \uline{partie fixe}
\[
T_{\handwritten{\ell}}(A^{\mathrm o})^{\mathrm f}\handwritten{\lhook\joinrel\longrightarrow}T_{\handwritten{\ell}}(A)
\]
est définie comme le pro-$\handwritten{\ell}$-groupe de Barsotti-Tate formé des $({}_{\handwritten{\ell^\nu}}A^{\mathrm o})^{\mathrm f}$. Nous allons maintenant donner une définition de la partie torique (2.3.2) de $T_{\handwritten{\ell}}(A^{\mathrm o})$, indépendante de l'hypothèse $\handwritten{\ell}\ne p$ faite dans 2.3, mais supposant en revanche $S$ \uline{complet} (*). Pour ceci, introduisons pour tout entier $n\geq0$ le schéma artinien local
\[
S_n=\operatorname{Spec}(V/\underline{\mathrm m}^{n+1})\ ,\tag*{(5.1.1)}
\]
$n$.ème voisinage infinitésimal de $s$ dans $S$, et pour tout $S$-schéma $X$ soit $X_n=X\times_S S_n$. Comme $(A^{\mathrm o})_n=(A_n)^{\mathrm o}$ est lisse sur $S_n$, il est connu (SGA 3 IX 3.6 bis) que $T_{\mathrm o}$ est la fibre spéciale d'un unique sous-tore
\[
T_n\handwritten{\subset}A_n^{\mathrm o}\ .\tag*{(5.1.2)}
\]
D'après l'unicité des $T_n$, ceux-ci se déduisent les uns des autres par les changements de base $S_{n'}\handwritten{\longrightarrow} S_n$. En d'autres termes, il définissent un \uline{sous-tore formel} $\hat T$ du groupe $\hat A^{\mathrm o}$ (dans la catégorie des schémas formels sur $S$) complété formel de $A^{\mathrm o}$ le long de la fibre spéciale :
\[
\hat T\handwritten{\subset}\hat A^{\mathrm o}\ .\tag*{(5.1.3)}
\]
\noindent\rule{6em}{.4pt}\par
(*) Cette hypothèse n'est d'ailleurs pas essentielle, comme on montre dans 6.1 ci-dessous.
% archive sheet number: 208
% pencil number at foot of upright typescript: 369
% stamp: numbered label 369
